\section{Method}
\label{sec:method}

\subsection{Setup}
\label{sec:method-setup}

Our setting is the standard mobile GUI agent loop. At each step the agent receives a screenshot of the current Android screen and a natural-language goal; it emits one of a fixed set of low-level actions (click on a coordinate or a labeled element, type a string, scroll, navigate, or signal task completion); the action is executed through the device's accessibility layer, producing the next screenshot. The episode ends when the agent emits a completion status, when an external grader verifies success, or when a per-task step budget is exhausted.

We use Qwen3-VL-8B-Instruct~\citep{qwen3vl2025} as the vision-language backbone and wrap it in the M3A two-phase loop introduced together with AndroidWorld~\citep{rawles2024androidworld}: an action-selection call conditions on the current screenshot plus a textual summary of the history, and a separate summary call distills the executed step into the history seen at the next action call. The action space and grading are the standard AndroidWorld schema; episode success is the rule-based grader's verdict, and self-reported \texttt{status:complete} alone does not count.

The action-selection prompt has a fixed prefix, an optional slot for the agent's in-context knowledge $F$, and a dynamic block containing the goal, history, and current UI elements. Different baselines instantiate $F$ differently---empty in the zero-shot baseline, a static $L_C$ text in the static-knowledge baseline, an evolving FSM under our method---but everything else is byte-identical across baselines, so comparisons isolate the contribution of $F$ rather than confounding it with prompt differences.

\subsection{Two-layer FSM}
\label{sec:method-fsm}

The core design choice in this paper is to structure $F$ as a finite state machine separated into two layers that differ in their transferability. This separation is the unit of what gets aggregated, retrieved, and mutated across apps; it is also the structural commitment that gives our system room to specialize without breaking transfer.

\textbf{LAYER~1 (app-specific).} A directed graph whose nodes are screens (e.g.\ \texttt{FILE\_LIST}, \texttt{NEW\_FILE\_DIALOG}) and whose edges are interactions that transition between them. Each node carries visual cues (e.g.\ \emph{``red plus FAB bottom-right''}), resource hints (the strings the model is likely to read on that screen), and known dead-ends. Every entry binds to one app's UI and is therefore non-transferable.

\textbf{LAYER~2 (category-generic).} A set of task-shape archetypes such as \texttt{ADD\_ENTRY}, \texttt{QUERY\_INFO}, \texttt{REMOVE\_ENTRY}, and \texttt{FILTER\_LIST}, each with four fields: \emph{precondition}, \emph{abstract\_steps}, \emph{failure\_modes}, and \emph{verification\_checklist}. LAYER~2 is written without any app-specific token. A linter (\texttt{lint\_layer2}) rejects any LAYER~2 block, or any \emph{diff} during online evolution, that contains an app name or a resource identifier from the source pool. This enforced separation is what makes the upper layer the structurally guaranteed unit of cross-app transfer.

\textbf{The category-level abstract library $L_C$.} Same-category source apps share their LAYER~2 content into one merged $L_C$ file per Play Store category. The merger keeps only claims supported by multiple sources, explicitly enumerates cross-source failure modes, and never names a specific app. At deployment, $\mathrm{resolve\_L\_C}(a_{\text{tgt}})$ returns the $L_C$ for the target app's category if that category is in the source pool, or $\bot$ otherwise. The latter case (Tier-C) gives our evaluation a built-in null control: with no $L_C$ to inject, the prompt degenerates to the zero-shot baseline byte-identically.

\textbf{How the FSMs are built.} For each source app we collect $K{=}5$ trajectories under the zero-shot baseline, concatenate them into a compressed text bundle, and ask Claude Opus 4.7~(1M context) to synthesize a two-layer FSM in JSON. The bundle fits in a single call and the model is prompted to keep LAYER~2 strictly app-agnostic; we then lint the result and audit it qualitatively. The 12 source FSMs feed the $L_C$ aggregator, which uses Claude Opus again on a per-category basis to merge same-category LAYER~2 blocks. The resulting six $L_C$ files (Productivity, Tools, Finance, Music \& Audio, Communication, Photography) cover the six categories in our source pool.

\subsection{The EvoFSM-RL adaptation loop}
\label{sec:method-loop}

At adaptation time we maintain a population $\mathcal{S}_a = \{(F_i, \mu_i, \sigma_i)\}$ of FSM variants for the current app $a$, each with a TrueSkill rating~\citep{herbrich2007trueskill}. One iteration of the loop proceeds as in Algorithm~\ref{alg:evofsmrl}.

\begin{table}[t]
\small
\centering
\begin{tabular}{p{0.42\textwidth}}
\hline
\textbf{Algorithm 1: One EvoFSM-RL iteration} \\
\hline
\textbf{Input:} population $\mathcal{S}_a$, task set $\mathcal{T}_a$, policy $\pi_\theta$, frozen reference LLM $\pi_{\text{ref}}$, GRPO buffer $\mathcal{B}$ \\
\hline
1. Sample task $t \sim \mathcal{T}_a$ \\
2. Select $M$ variants $\{F_i\}_{i=1}^M$ from the latest-$K$ window of $\mathcal{S}_a$ with $p_i \propto \exp((\mu_i + \lambda\sigma_i)/T)$ \\
3. For each $i$, run $N$ rollouts of $\pi_\theta$ under prompt $F_i$ on task $t$ with paired seeds; record reward $r_{i,j}$ and per-step log-probs \\
4. Update $\{(\mu_i, \sigma_i)\}$ by TrueSkill from the ranking of $\{V_i = \tfrac{1}{N}\sum_j r_{i,j}\}$ \\
5. Append the $M{\times}N$ trajectories to $\mathcal{B}$; if $|\mathcal{B}| \ge K_{\text{buf}}$ run a GRPO step on $\mathcal{B}$ and clear it (Section~\ref{sec:method-grpo}) \\
6. Let $k = \arg\max_i V_i$. If $r_k > 0$ or $i \bmod 3 = 0$: synthesize a LAYER~2 diff with $\pi_{\text{ref}}$, apply it to $F_k$, append the child to $\mathcal{S}_a$ with $\sigma_k + \Delta\sigma$ \\
\hline
\end{tabular}
\caption{One iteration of the EvoFSM-RL adaptation loop. The same procedure runs in Phase~1 (source-pool pretraining) and in Phase~3 (TTA on a target app). The only differences are which $\mathcal{T}_a$ we draw from and whether $F$ is allowed to mutate.}
\label{alg:evofsmrl}
\end{table}

Three design choices in this loop are load-bearing.

\textbf{Mutation uses a frozen reference, not the trainable policy.} We prompt Claude Opus 4.7 with a layered-reflection template over the recent trajectory window and ask for a JSON-typed diff over LAYER~2 only. The lint pass guarantees no leakage. We deliberately do not use $\pi_\theta$ itself for the reflection because, in joint training, the RL-updating policy drifts as it learns, and using a moving target as a critic of its own behavior degrades the quality of the proposed edits---a finding ablated in~\citet{zhang2026espl}.

\textbf{Paired seeds within a tournament round.} All $M$ variants are rolled out on the same task and the same rollout seed within an iteration. The only source of reward difference between variants is therefore the FSM content, which lets TrueSkill compare them on identical content rather than confounding the comparison with task-difficulty noise.

\textbf{LAYER~1 is frozen during online adaptation.} LAYER~1 is built offline from source-pool trajectories and is not mutated during B3 or B4. Online adaptation modifies only LAYER~2. This is what keeps online evolution bounded to the transferable part of the knowledge; if LAYER~1 were also mutable, a few iterations on a target app could absorb app-specific edits into the wrong layer and silently break transfer downstream.

\subsection{Group-relative policy optimization}
\label{sec:method-grpo}

The weight update is a group-relative policy gradient~\citep{shao2024deepseekmath} computed on the buffered trajectories. We compute the reward for each rollout as
\begin{equation}
R(t,\tau) = s + w_e \cdot \max\!\left(0,\, 1 - \tfrac{n_{\text{steps}}(\tau)}{B(t)}\right),
\label{eq:reward}
\end{equation}
where $s \in \{0, 0.5, 1\}$ is the AndroidWorld grader's verdict, $B(t)$ is the per-template step budget, and $w_e = 0.2$ keeps the efficiency bonus small enough that success dominates.

Following the single-task M-prompts-by-N-rollouts shape of Algorithm~1 in~\citet{zhang2026espl}, we key groups by the $(F_i, t)$ pair: within each iteration, the $N$ rollouts under variant $F_i$ on task $t$ form one group, and the advantage of rollout $j$ is
\begin{equation}
A_{i,j} = r_{i,j} - V_i, \qquad V_i = \tfrac{1}{N}\sum_{k=1}^{N} r_{i,k}.
\label{eq:advantage}
\end{equation}
The signal $A_{i,j}$ is purely a within-(FSM, task) deviation from the local mean. A group of size~1 has $A \equiv 0$ and contributes nothing to the update, which is why we require $N \ge 2$ at all times.

The LoRA-only gradient on a buffer of active trajectories is
\begin{equation}
\nabla_\theta \mathcal{L} = -\sum_{j \in \mathcal{B}_{\text{active}}} \frac{A_j}{T_j} \sum_{t=1}^{T_j} \nabla_\theta \log \pi_\theta(a_t \mid s_t, F),
\label{eq:grpo-loss}
\end{equation}
where $T_j$ is the number of agent steps in rollout $j$. The $1/T_j$ normalization keeps long and short rollouts on the same gradient scale; without it, the raw gradient norm scales with trajectory length and gets clipped uniformly by $\|\nabla\|_{\max}$, eliminating the actual learning direction in favor of trajectory-length artifact. We freeze the base vision-language model and update only LoRA matrices on the attention projections $\{q_{\text{proj}}, v_{\text{proj}}\}$~\citep{hu2021lora}; KL regularization is implicit in the LoRA-only formulation and we set $\beta = 0$.

\subsection{Two-phase application}
\label{sec:method-phases}

The same EvoFSM-RL loop runs at two different points in the pipeline. The decision to use the \emph{same} loop in both places---rather than train one model offline and fine-tune another online---is what lets us call the deployment-time procedure a well-defined operation rather than an ad-hoc fine-tune.

\textbf{Phase~1: shared-LoRA pretraining.} A single training run over the source pool produces a shared LoRA adapter $\pi^{\text{pre}}_\theta$. Each iteration samples one source app uniformly, samples one of its tasks, runs $N$ rollouts under the app's frozen static FSM, and accumulates the trajectories into a single GRPO buffer. The buffer's groups are still keyed by $(F_i, t)$, so gradients from rollouts on different source apps flow through their respective task groups into the same adapter. We freeze FSM mutation in Phase~1 (a deliberate scope reduction---the static FSMs and $L_C$ are already produced by simpler non-RL pipelines as inputs to this phase).

\textbf{Phase~3: deployment-time adaptation.} For each target app we run the same loop on the app's $T_{\text{adapt}}$ split, starting from $\pi^{\text{pre}}_\theta$ and seeding the population from the category-matched $L_C$. Both the FSM population and the LoRA adapter evolve online. After a fixed budget of adaptation iterations (20 in our experiments), the LoRA adapter and the FSM champion are frozen, and the agent is evaluated on the held-out $T_{\text{eval}}$ split. This is the test-time half of the framework, and what we call B4 in the ablation.

A small comment on what Phase~2 looks like in our implementation. The full algorithmic description specifies a target-app cold-start that uses $\pi^{\text{pre}}_\theta$ to roll out a few exploration trajectories on the target app and synthesize a target-app LAYER 1 before adaptation begins. We collapse this to a single category look-up: $L_C \leftarrow \mathrm{resolve\_L\_C}(a_{\text{tgt}})$. The target app's LAYER 1 is rediscovered in-context during the B3/B4 loop itself, rather than offline. This works because the visual encoder is strong; we discuss when it fails in Section~\ref{sec:limitations}.

The Phase~1 $\to$ Phase~3 split is what makes the joint paradigm tractable at the deployment-time budget. Twenty iterations from a randomly initialized LoRA are not enough to learn anything meaningful, whereas twenty iterations starting from $\pi^{\text{pre}}_\theta$ exercise the policy in a regime where the local gradient signal is small but the starting point already encodes a useful prior.
